\begin{proof}
\begin{enumerate}[label=\arabic*.]
\item For each proposed pair, \ref{maps:src:F.2.9} gives an
irreducible rational rank-five quotient $W$.  Its cusp has Jordan
type $(2,2,1)$.  One finite polynomial is
$(T-1)\Phi_3(T)^2$ and the other is $(T-1)\Phi_d(T)$ for
$d=5,8,10$.  For the two additional assertions, replace
$\Phi_3^2$ by $\Phi_4^2$ or $\Phi_6^2$, with $d=5$.
In all these cases a commuting endomorphism of $W$ restricts to a
rational scalar on the one-dimensional fixed space of the second
generator.  This restriction is injective: a nonzero endomorphism
of an irreducible representation is invertible and cannot vanish
on that line.  Consequently
\[
 \operatorname{End}_{\Gamma}(W)=\mathbb Q.
\]

\item Consider $\mathcal H=\operatorname{Hom}(W,W^*)$.
The local fixed dimensions are $9,5,13$.  At the repeated finite
site, pairing reciprocal eigenspaces of dimensions $1,2,2$ gives
$1+4+4=9$.  At the distinct finite site it gives five.
For the cusp, the ordered pairs of blocks of sizes $2,2,1$
contribute the sum of their minimum sizes, namely thirteen.
The invariant-sheaf Euler formula therefore gives
\[
 \chi(\mathbb P^1,j_*\mathcal H)=-25+9+5+13=2.
 \tag{F.31}\label{maps:eq:F.31}
\]
Here degree zero is $\operatorname{Hom}_{\Gamma}(W,W^*)$,
and degree two is dual to
$\operatorname{Hom}_{\Gamma}(W^*,W)$.  At least one is nonzero,
so irreducibility gives a nondegenerate invariant bilinear form
$J:W\longrightarrow W^*$, inverting the map if necessary.
Its alternating part is equivariant; if nonzero it would be
invertible, which is impossible in odd dimension.  Thus $J$ is
symmetric.

\item An invariant bivector is an alternating equivariant map
$K:W^*\longrightarrow W$.  The actual comparison map
\[
 K\longmapsto KJ\in\operatorname{End}_{\Gamma}(W)
\]
makes it a scalar, so $K$ is a scalar multiple of the symmetric
map $J^{-1}$.  Hence $K=0$.  Applying the same argument to an
alternating map $W\longrightarrow W^*$ gives
\[
 (\Lambda^2W)^\Gamma=(\Lambda^2W^*)^\Gamma=0.
\]
Both degree zero and degree two of $j_*\Lambda^2W$ vanish.
Its Euler characteristic must therefore be nonpositive.

\item The repeated finite spectrum has four fixed bivectors and
the distinct spectrum has two.  At the cusp, write the unipotent
module as two blocks of size two and one block of size one.
Their exterior square has two trivial blocks from the individual
wedges, blocks of sizes three and one from the product of the
two size-two blocks, and two further size-two blocks.  Its fixed
rank is six.  Consequently
\[
 \chi(\mathbb P^1,j_*\Lambda^2W)=-10+4+2+6=2,
\]
contradicting the preceding step.  This proves all five stated
exclusions.  Removing the three excluded order pairs from
\eqref{maps:eq:F.28} leaves exactly \eqref{maps:eq:F.30}.
\end{enumerate}
\end{proof}
